\title{Byte Latent Transformer: Patches Scale Better Than Tokens}

\begin{abstract}
We present the Byte Latent Transformer ({\textsc BLT}), a novel large language model (LLM) architecture operating at the byte level that achieves, for the first time, performance comparable to tokenization-based LLMs at scale, while delivering notable enhancements in inference speed and resilience. 
{\textsc BLT} processes raw bytes by aggregating them into patches of variable length, which act as the principal computational elements. The model determines patch boundaries by assessing the entropy of the next byte, directing greater computational and model resources to areas with higher informational complexity. We introduce the first flop-controlled scaling analysis for byte-level models, exploring sizes up to 8B parameters and 4T training bytes. Our findings validate that it is possible to scale models directly from unprocessed byte sequences without relying on predefined vocabularies. Efficiency in both training and inference is achieved by adaptively forming larger patches in segments where byte sequences are more predictable; additionally, we observe qualitative gains in reasoning capabilities and generalization to rare or complex inputs. All in all, at a constant inference computational budget, {\textsc BLT} demonstrates substantially improved scaling over conventional token-based architectures by jointly increasing patch lengths and model capacity.
\end{abstract}

\section{Introduction}
\label{section:intro}

We present the Byte Latent Transformer~(\textbf{BLT}), an architecture that entirely dispenses with tokenization by learning directly from sequences of raw bytes. For the first time, BLT achieves parity with tokenization-dependent models at large scales, while offering marked gains in both computational efficiency and robustness (\S\ref{sec:robustness}). While large language models (LLMs) are predominantly trained using end-to-end objectives, the tokenization phase remains an ad hoc pre-processing mechanism that segments byte sequences into a predefined token vocabulary. This design imposes inductive biases in string compression and introduces several weaknesses, including heightened sensitivity to domains and modalities~\citep{dagangetting}, vulnerability to input noise~(\S\ref{sec:robustness}), limitations in capturing orthographic structure~\citep{edman2024cute}, and disproportionate handling of multiple languages~\citep{liang2023xlm,petrov2024language,limisiewicz2024myte}.

Historically, tokenization has played a critical role, as training llms directly on byte-level inputs incurs excessive computational expense at large scales owing to the increased sequence lengths~\citep{xue2022byt5}. To address this, earlier research has focused on leveraging more computationally efficient self-attention models~\citep{el2020characterbert, clark2022canine} or adopting attention-free approaches~\citep{wang2024mambabyte}~(\S\ref{section:related_work}). These strategies, however, mainly yield benefits when training \textit{small models}. When scaling up, the primary computational burden within Transformer architectures shifts to the substantial feed-forward network layers that must process each byte, rendering the cost associated with the attention mechanism relatively minor.

In order to optimize compute distribution, we introduce a flexible, adaptive strategy for assembling bytes into \textit{patches}~(\S\ref{section:patching}), coupled with a novel model framework that integrates both byte-level and patch-level information. Distinct from standard tokenization, {\textsc BLT} dispenses with a predetermined patch vocabulary. Instead, it uses compact, learnable encoder and decoder components to transform any collection of bytes into latent patch embeddings. Our findings demonstrate that this yields a \textit{more} efficient use of computational resources than approaches relying on tokenization.

Tokenization-based language models assign identical computational resources to each token. This design sacrifices efficiency for performance because the tokenization process relies on compression strategies that do not necessarily reflect the actual difficulty of making predictions. A core principle of our model is to enable dynamic computation allocation that matches the demands of different prediction tasks. For instance, generating the last few characters of a word often requires less computational effort, as these represent straightforward, low-entropy predictions—unlike selecting the initial word of a sentence, which can be much more complex. This principle is operationalized in {\textsc BLT}'s architecture~(\S\ref{section:architecture}), which includes three distinct transformer modules: a pair of lightweight byte-level \textit{local models} and a robust global \textit{latent transformer}~(\autoref{fig:arch_high_level}). 

To adaptively form byte groupings, and thus manage computational focus, {\textsc BLT} utilizes next-byte prediction entropy as a basis for data segmentation. This process yields contextual clusters of bytes that maintain consistent information density.

We introduce the first comprehensive scaling analysis of byte-level models regulated by floating-point operations (flops), extending to 8B parameters and 4T bytes of training data. This demonstrates that it is feasible to train large-scale models directly on raw bytes, without relying on predetermined vocabulary tokenization. Notably, {\textsc BLT} achieves parity with Llama 3 in flop-controlled training performance\footnote{Model computational expenditure is quantified as the total number of Floating Point OPerations (flops) required.}, while reducing inference flops by up to 50\%~(\S\ref{section:scaling}).

Moreover, we find that operating directly on raw bytes substantially enhances the model's ability to capture rare and diverse phenomena in the data distribution. {\textsc BLT} exhibits improved robustness to noisy data compared to token-based models and demonstrates superior capabilities in character-level comprehension tasks, including orthographic representation, phonological analysis, and performance in low-resource machine translation~(\S\ref{sec:robustness}).

Furthermore, the {\textsc BLT} architecture permits concurrent scaling of both model size and patch size without exceeding a fixed inference flop budget. Using larger patch sizes reduces computation per inference step, which enables redirecting saved compute resources to increase the capacity of the global latent transformer, as it is invoked less frequently. Our inference flop-constrained scaling experiments~(\autoref{fig:fixed_inference_scaling}) reveal considerably improved scaling behavior relative to models utilizing fixed-vocabulary tokenization.

To conclude, this work provides the following advancements: 1) We present {\textsc BLT}, a byte latent llm model that allocates computational resources adaptively to increase flop efficiency, 2) We demonstrate that our approach matches Llama 3 in training flop parity up to the 8B parameter scale, while also enabling significant trade-offs—accepting marginal declines in evaluation scores to attain up to 50\% improvements in flop efficiency, 3) The {\textsc BLT} framework introduces a novel scalability axis for llms, allowing enlargement of model capacity without exceeding a constant inference budget, 4) Our results show that {\textsc BLT} models possess heightened resilience to input perturbations and recognize sub-word information often overlooked by token-centric llms. All training and inference resources for {\textsc BLT} are publicly available at~\url{https://github.com/facebookresearch/blt}.

\section{Patching: From Individual Bytes to Groups of Bytes}
\label{section:patching}

Dividing bytes into \textit{patches} enables {\textsc BLT} to flexibly distribute computational resources according to the surrounding context. Several approaches to partitioning bytes into patches are illustrated in Figure~\ref{fig:patching_types}. More formally, a patching function $f_p$ takes an input sequence of bytes $\pmb{x}=\{x_i\,|\,i=1,\ldots, n\}$ of length $n$ and groups it into a series of $m < n$ patches, denoted as $\pmb{p}=\{p_j\,|\,j=1,\ldots,m\}$. The segmentation occurs via mapping each $x_i$ to \{0,1\}, where a value of 1 designates the initiation of a new patch. The computational load associated with processing data, for both token-based and patch-oriented models, depends largely on the number of main Transformer operations. For {\textsc BLT}, this workload is dictated by the number of patches needed to represent the sequence, as defined by the selected patching function. As a result, the mean number of bytes per patch—referred to as the \textit{patch size}—serves as a key metric for estimating computational expense during both model training and inference, given a particular patching strategy~(\S\ref{section:flops}). In the following sections, we detail three distinct patching functions: partitioning by a constant byte count per patch~(\S\ref{section:static-patch}), segmenting on whitespace boundaries~(\S\ref{section:white-patch}), and adaptively patching guided by the local entropy calculated from a compact byte-level language model~(\S\ref{section:dyn-patch}). We conclude with a discussion on incremental patching and how the concept of tokenization diverges from that of patching~(\S\ref{section:bpe-patch}).

\subsection{Strided Patching Every K Bytes}
\label{section:static-patch}

A simple and commonly used method for grouping bytes is to segment them into uniform patches of length $k$, as demonstrated in MegaByte~\citep{yu2023megabyte}. This technique, which employs a consistent stride, facilitates straightforward implementation during both training and inference phases, and conveniently allows for adjustment of the average patch length. As a result, it also provides a direct way to manage computational costs. Despite these advantages, fixed-size patching presents notable limitations. Primarily, computational resources are not adaptively focused where they are most required: a model step $j$ could be underutilized by allocating resources to predict trivial elements such as whitespace in code, or, conversely, could lack sufficient capacity to process content-rich bytes like those found in mathematical expressions. Additionally, this strategy produces irregular and context-insensitive segmentations of identical byte sequences, such as cases where an identical word may be partitioned in disparate ways.

\subsection{Space Patching}
\label{section:white-patch}

\citet{slagle2024spacebyte} introduces a straightforward but impactful modification to strided patching, in which patches are generated whenever a space-like byte\footnote{
    Space-like bytes refer to any byte excluding Latin characters, digits, or UTF-8 continuation bytes; furthermore, every patch is required to include at least one byte that is not space-like.
} is encountered—these typically correspond to natural linguistic boundaries in numerous languages. The method, called Space patching, assigns a latent transformer computation step (i.e., additional floating point operations) for each word, thereby ensuring a consistent patching scheme for words across different sequences, and dedicating computational resources specifically to challenging predictions that frequently occur following space-like bytes. Consider, for instance, answering the question ``Who composed the Magic Flute?\uline{\hspace{.6em}}''. Generating the first byte of the answer is substantially more challenging than generating subsequent bytes once the initial ``M'' is provided, as the initial character dramatically narrows the pool of plausible options—rendering the prediction of the rest of ``Mozart'' relatively straightforward. Nevertheless, space patching is not universally suitable for all languages or application domains, and, crucially, does not offer flexibility in patch size. Subsequently, we present an alternative patching approach, which builds on the observation that predicting the initial byte in words typically poses the greatest challenge, while offering an intuitive solution for managing patch size.

\subsection{Entropy Patching: Using Next-Byte Entropies from a Small Byte LM}
\label{section:dyn-patch}

Instead of employing traditional rule-based heuristics like whitespace, our method utilizes a data-driven strategy to detect regions where the model exhibits high uncertainty in predicting the next byte. We propose \textit{entropy patching}, a technique that determines patch boundaries by leveraging entropy estimates.

We train a small byte-level auto-regressive language model on the training data for {\textsc BLT} and compute next byte entropies under the LM distribution $p_{e}$ over the byte vocabulary $\mathcal{V}$:

\begin{equation}
H(x_i) = \sum_{v \in \mathcal{V}} p_{e}(x_i=v|\pmb{x}_{<i}) \log p_{e}(x_i=v|\pmb{x}_{<i})
\end{equation}

We utilize two approaches to detect patch boundaries based on the calculated entropies $H(x_i)$. The first approach involves selecting points where the entropy surpasses a predefined global threshold, as depicted in~\autoref{fig:patching}. The alternative method detects points exhibiting higher entropy relative to the immediately preceding value. This latter strategy may also be viewed as marking positions that interrupt a roughly monotonic decline in entropy throughout the patch.

\begin{align*}
    \text{Global Constraint}\hspace{1cm}H(x_t)& > \theta_g\\
    \text{Approx. Monotonic Constraint}\hspace{1cm}H(x_t)& - H(x_{t-1}) > \theta_r
\end{align*}

Patch boundaries are determined through a simple preprocessing procedure applied at the dataloading stage. This approach contrasts with that of~\citet{nawrot-etal-2023-efficient}, in which a classifier is used to infer patch boundaries based on entropy measures. In our study~(\S\ref{section:experiments}), we evaluate both techniques for separating low-entropy from high-entropy bytes.

\subsection{The Byte-Pair Encoding (BPE) Tokenizer and Incremental Patching}
\label{section:incremental-patching}
\label{section:bpe-patch}

A large number of contemporary llms, such as our baseline Llama 3, employ a subword tokenizer like bpe~\citep{gage1994new, sennrich-etal-2016-neural}. Throughout this work, ``tokens'' denote byte groupings selected from a \textit{finite} vocabulary defined before training begins, whereas ``patches'' refer to adaptively formed sequences that do not rely on a pre-established vocabulary. The key distinction between tokens and patches is that with tokens, the underlying byte information is not directly accessible to the model.

An important advancement offered by {\textsc BLT} compared to tokenization-based architectures is its reconfiguration of the balance between vocabulary size and computational requirements. Traditional llms face a scenario where expanding the vocabulary results in average tokens encompassing more bytes, which reduces the number of steps required during inference. However, this simultaneously expands the dimensionality of the model’s final projection layer, increasing computational burden. This compromise limits how much token-based methods can adjust token length and inference expenditure. As an illustrative case, Llama 3 managed to raise the average token length from 3.7 to 4.4 bytes, but this came at the expense of making its embedding table four times larger than that of Llama 2.

During generation, {\textsc BLT} must determine, at each step in the byte sequence, whether it has reached a patch boundary—a choice that controls if additional computation is triggered through the Latent Transformer. This assessment must be performed without knowledge of the subsequent, as-yet-unproduced bytes in the sequence. Consequently, the process of patching cannot rely on the content of future bytes for segmenting the sequence. More precisely, a patching method $f_p$ fulfills the criterion of incremental patching if:
$$f_p(\pmb{x}_{<i}) = f_p(\pmb{x})_{<i}$$
In contrast, bpe does not qualify as an incremental patching scheme, because the way a prefix is tokenized can vary depending on how the sequence continues, thus failing to uphold the above property\footnote{Introducing a dedicated delimiter token to mark patch boundaries allows bpe to function as an incremental patching scheme, albeit at the cost of elongating the byte sequence.}.

\section{{\textsc BLT} Architecture}
\label{section:architecture}
The {\textsc BLT} system consists of a main global autoregressive language model designed to process patch-level representations. In addition, it employs a pair of lightweight local models: one responsible for encoding byte sequences into patches and another for converting patch representations back into bytes~(\autoref{fig:arch_high_level}).

\subsection{Latent Global Transformer Model}

\textit{The Latent Global Transformer} functions as an autoregressive transformer model $\mathcal{G}$ comprised of $l_{\mathcal{G}}$ layers, which processes sequences of latent input patch representations $p_j$ and produces corresponding output patch representations $o_j$. In this work, the notation $j$ consistently refers to patches, while $i$ indexes individual bytes. The global transformer adopts a block-causal attention masking strategy~\citep{dubey2024llama}, which enforces that each position in the sequence can attend only to patches up to and including the current patch within a document. Because the global model constitutes the most computationally intensive component during both pre-training and inference, selectively determining its application enables dynamic adjustment of computational resources depending on the complexity of specific input or output regions.

\subsection{Local Encoder}
\label{section:local}

\textit{The Local Encoder Model}, represented as $\mathcal{E}$, is a compact model grounded in the transformer architecture, consisting of $l_{\mathcal{E}}$ layers, where $l_{\mathcal{E}} << l_{\mathcal{G}}$. Its central objective is to rapidly transform a sequence of input bytes $b_i$ into informative patch-level representations $p_j$. One notable modification from standard transformers is the integration of a cross-attention layer subsequent to every transformer layer, which serves to aggregate byte-level embeddings into their respective patch representations~(\autoref{fig:crossattn}). 

The process begins by embedding each input byte $b_i$ through a matrix in $\mathbb{R}^{256 \times h_{\mathcal{E}}}$, yielding vectors $x_i$. Optionally, these base embeddings can be enriched with supplementary features via hash-embeddings~(\S\ref{sec:hash-emb}).

This sequence of embeddings is then passed through a stack of transformer layers interleaved with cross-attention layers~(\S\ref{sec:enc-cross-attn}), which iteratively refine the representations into patch-level vectors $p_i$. These output patches are intended for consumption by the global transformer module $\mathcal{G}$.

The attention mechanism within the transformer blocks employs a \textit{local block causal} attention pattern: each byte considers information from at most $w_{\mathcal{E}}$ preceding bytes, potentially crossing current patch boundaries, but always respecting document segmentation. The next sections provide a deeper exploration of the embedding strategies and the specifics of the cross-attention architecture.

\subsubsection{Encoder Hash n-gram Embeddings}
\label{sec:ngram-emb}
\label{sec:hash-emb}
An essential aspect of generating resilient and expressive representations at each position $i$ involves capturing context from preceding bytes. In {\textsc BLT}, this is accomplished by representing the byte $b_i$ on its own as well as within the context of a byte n-gram. Specifically, at every step $i$, we begin by forming byte-grams

\begin{align}
    g_{i,n}=\{b_{i-n + 1},\ldots, b_i\}
\end{align}

for each byte position $i$ and for each $n$ ranging from three to eight.\footnote{
We exclude byte-grams of length $n$ or greater when $i < n$. }

Next, we propose hash $n$-gram embeddings, which utilize a hash function to project every possible byte $n$-gram onto a corresponding index in an embedding matrix $E_{n}^{hash}$ of constant size, for every $n$ in $\{3, 4, 5, 6, 7, 8\}$~\citep{bai2010learning}.
The associated embedding from this table is combined additively with the base byte embedding, then the sum undergoes normalization prior to being fed into the local encoder. The final, enriched embedding is expressed as

\begin{align}
e_i &= x_i + \sum_{n = 3,...,8}E_{n}^{hash}(\text{Hash}(g_{i,n})) \\
\text{where, Hash}(g_{i,n}) &= \text{RollPolyHash}(g_{i,n}) \% |E_{n}^{hash}|
\end{align}

The value $e_i$ is normalized by dividing by the total number of $n$-gram sizes, incremented by one, and the RollPolyHash function—specified in Appendix~\ref{appendix:rollpolyhash}—is employed. Section~\ref{section:ablations} discusses an ablation study examining how varying $n$-gram hash embedding parameters, specifically the choice of $n$ and the size of the embedding table, affect flop-controlled scaling law patterns. Alongside hash-based $n$-gram embeddings, experiments utilizing frequency-based $n$-gram embeddings were conducted; further information regarding these experiments can be found in Appendix~\ref{appendix:ngrams}.

\subsubsection{Encoder Multi-Headed Cross-Attention}
\label{sec:enc-cross-attn}
Our approach largely mirrors the input cross-attention mechanism described in the Perceiver architecture~\citep{jaegle2021perceiver}. However, instead of employing a fixed set of latent variables, our latent representations directly map to dynamic patch representations~(\autoref{fig:crossattn}), and the attention mechanism is restricted to only the bytes belonging to their respective patches. Specifically, for each patch $p_j$, we generate a query vector by aggregating the byte-level representations within $p_j$, which is then subjected to a linear transformation via $\mathcal{E}_{C} \in \mathbb{R}^{h_{\mathcal{E}} \times (h_{\mathcal{E}} \times U_{\mathcal{E}})}$, where $U_{\mathcal{E}}$ signifies the count of encoder cross-attention heads. Explicitly, let $f_{\text{bytes}}(p_j)$ be the sequence of bytes for patch $p_j$, we then compute

\begin{align}
P_{0,j} &= \mathcal{E}_{C}(f_\text{bytes}((p_j)), f~ \text{is a pooling function} \\
P_l &= P_{l-1} + W_o\left(\text{softmax}\left(\frac{QK^T}{\sqrt{d_k}}\right)V\right) \\
\text{where } Q_j &= W_q(P_{l-1,j}), K_i = W_k(h_{l-1,i}), V_i = W_v(h_{l-1,i}) \\
h_l &= \text{Encoder-Transformer-Layer}_l(h_{l-1})
\end{align}

Here, $P \in \mathbb{R}^{n_p \times h_{\mathcal{G}}}$ denotes the representations for $n_p$ patches that are input to the global model. To form $P$, we pool the byte embeddings $e_i$ associated with each individual patch $p_j$ at initialization. The matrices $W_q$, $W_k$, $W_v$, and $W_o$ are the learnable projection parameters for constructing the query, key, value, and output representations, respectively. Keys and values are produced by projecting the byte representations $h_i$ from the preceding layer, or from $e_i$ if it is the first layer. To ensure each query $Q_j$ only aggregates information from its own patch, we implement a mask such that attention weights are applied exclusively to the bytes within the same patch $j$. The multi-head attention mechanism operates over the $Q, K, V$ representations. Noting that patch representations typically have a greater dimensionality ($h_{\mathcal{G}}$) compared to $h_{\mathcal{E}}$, we represent $P_l$ as a set of heads each of size $h_{\mathcal{E}}$ in the cross-attention stage, subsequently concatenating them to yield a vector of size $h_{\mathcal{G}}$. Furthermore, a pre-LayerNorm is applied to queries, keys, and values prior to attention; positional encodings are omitted within this cross-attention module. A residual connection wraps around the cross-attention block to enhance information flow.

\subsection{Local Decoder} 

The local decoder $\mathcal{D}$, much like the local encoder, is a transformer-based architecture designed to be lightweight, employing $l_{\mathcal{D}} << l_{\mathcal{G}}$ layers. Its role is to transform the sequence of global patch embeddings $o_j$ into output raw bytes, $y_i$. Functioning autoregressively, the local decoder generates each byte by conditioning on the sequence of previously output bytes, leveraging the internal representations generated by the local encoder for that byte sequence. Its architecture alternates $l_{\mathcal{D}}$ layers that combine cross-attention mechanisms with transformer layers: the cross-attention is applied initially to map the patch representations into the byte space, after which the transformer's operations refine the sequence at the byte level.

\subsubsection{Decoder Multi-headed Cross-Attention} 

In the decoder cross-attention, the roles of the queries and key/values are interchanged i.e. the byte-representations are now the queries, and the patch representations are now the key/values. The initial byte-representations for the cross-attention are initialized as the byte embeddings from the last encoder layer i.e. $h_{l_{\mathcal{E}}}$. The subsequent byte-representations for layer $l$, $d_{l,i}$ are computed as:

\begin{align}
D_0 &= h_{l_{\mathcal{E}}} \\
B_l &= D_{l-1} + W_o\left(\text{softmax}\left(\frac{QK^T}{\sqrt{d_k}}\right)V\right), \\
  &\text{where } Q_i = W_q(d_{l-1,i}), K_i = W_k(\mathcal{D}_C(o_j)), V_i = W_v(\mathcal{D}_C(o_j)) \\
  D_l &= \text{Decoder-Transformer-layer}_l(B_l)
\end{align}

In this context, $W_k$ and $W_v$ are matrices used to project keys and values, which operate on a combination of a linear transformation and the split operation $\mathcal{D}_C$ applied to the concluding patch representations $o_j$ produced by the global model. The matrix $W_q$ serves as the projection for queries and acts on byte representations $d_{l-1}$ from the preceding decoder transformer layer, or on $h_{l_{\mathcal{E}}}$ for the initial layer. $W_o$ functions as the output projection matrix. As a result, $B$ has dimensions in $\mathbb{R}^{h_{\mathcal{D}} \times n_b}$, where $n_b$ denotes the total output byte count. The subsequent decoder layer representations $D_l$ are determined by processing the outcome of the cross-attention block $B$ through a decoder transformer layer. Analogous to the procedure in the local encoder cross-attention, multiple attention heads are utilized, pre LayerNorms are applied, positional embeddings are omitted, and the cross-attention component is surrounded by a residual connection.

\section{Experimental Setup}
\label{section:experiments}
We systematically design a suite of controlled experiments to ensure a rigorous comparison between {\textsc BLT} and tokenization-based models, with an emphasis on preventing {\textsc BLT} from gaining undue benefit through longer sequence contexts.

\subsection{Pre-training Datasets} In our experiments, all model sizes are pre-trained using two primary datasets: 1) The Llama 2 dataset~\citep{touvron2023llama}, consisting of 2 trillion tokens sourced from an assortment of publicly accessible materials. This dataset undergoes extensive cleaning and filtering procedures aimed at enhancing overall data quality. 2) {\textsc BLT}-1T: A newly curated collection amounting to 1 trillion tokens, also derived from a range of public datasets, which incorporates a portion of the pre-training corpus made available through Datacomp-LM~\citep{li2024datacomp}. The Llama 2 dataset serves as the foundation for scaling law analyses, following the guidelines established in~\cite{dubey2024llama} to identify optimal token counts and inform architecture decisions for {\textsc BLT}. Conversely, {\textsc BLT}-1T is employed for exhaustive pre-training to facilitate direct comparisons with Llama 3's performance on various downstream benchmarks. It should be noted that neither dataset contains material sourced from Meta's proprietary products or services. Additionally, for baseline comparisons involving models that rely on tokenization, we employ the Llama 3 tokenizer, which features a 128K token vocabulary; this tokenizer demonstrated stronger baseline results relative to its Llama 2 counterpart in our empirical evaluations.

\subsection{Entropy Model} In our study, the entropy model is implemented as a byte-level language model and is trained on the same data distribution as the full {\textsc BLT} model. By default, this model adopts a transformer architecture featuring 100M parameters, comprising 14 layers with a hidden size of 512, and employs sliding window attention over a span of 512 bytes. All other hyperparameter settings align with those used in our local and global transformer architectures. We also conducted experiments varying model capacities, receptive field sizes, and architectural choices, as described in \autoref{section:ablations}. Notably, for cases with sufficiently small receptive fields, the resulting trained entropy model can be compactly represented as an efficient lookup table.

\subsection{Entropy Threshold and Equalizing Context Length} In models employing entropy-based patching, we determine a patching threshold that results in a target average \textit{patch size} across the pretraining dataset mixture. In {\textsc BLT}, the \textit{patch size} is not constrained by the structure of tokenization, allowing it to be set freely, which in turn substantially affects the model’s utilized context length. To ensure a fair comparison and prevent larger patch sizes from benefitting from longer effective contexts, we keep the total byte count per batch fixed on average. Consequently, for models utilizing greater patch sizes, the sequence length is proportionally decreased. Specifically, with Llama 2 data, we adopt an 8k byte context window, whereas on the {\textsc BLT}-1T dataset, the average context window is set to 16k bytes, all the while keeping the mean batch size at 16M bytes.

Although the average batch size remains unchanged, dynamic patching techniques result in varying proportions between the number of bytes and patches per batch. To enhance efficiency, our implementation of {\textsc BLT} training aggregates patches into batches in a manner that eliminates the need for padding in the costly latent transformer, thereby maintaining a uniform patch count across all batches. During the training phase, input byte sequences are padded and, if necessary, truncated to lengths of 12k for Llama 2 and 24k bytes for the {\textsc BLT}-1T datasets. This strategy helps mitigate the risk of memory overflow caused by sequences containing exceptionally large patches.

\subsection{Entropy Model Context}
\label{section:entropy-context}
Our observations indicate that applying entropy patching results in increasingly larger patches when dealing with structured content such as multiple choice problems (refer to patching in an MMLU case shown in \autoref{fig:mmlu_patching}), which typically contain significant repetition. This pattern arises because repeated material in the entropy model context exhibits lower entropy. Consequently, for the comprehensive evaluation of {\textsc BLT}-Entropy with a patch size of 4.5, we clear the entropy context by introducing new lines and employ an approximate monotonicity constraint, as this approach proves more resistant to "entropy drift" linked to fluctuations in context length. This adjustment pertains solely to the way entropies are determined; we continue utilizing the same methodology for selecting the appropriate entropy threshold.

\subsection{FLOPs Estimation} 
\label{section:flops}

Our methodology for calculating transformer flops is primarily based on the formulas provided by Chinchilla~\citep{hoffmann2022training}, which include the flop counts associated with the feed-forward modules, qkvo projections in self-attention, as well as attention computation and output projection steps. Unlike their approach, we treat the input embedding stage as a highly efficient lookup table rather than as a dense matrix multiplication, effectively assigning it zero flops. In alignment with established literature, we approximate the backward pass to require double the flops of the forward computation.

To compute flops \textit{per byte} for {\textsc BLT} models, we add up the flops for the local encoder transformer, the global latent transformer, and the local decoder transformer, together with the cross attention blocks in the encoder and the decoder:

\begin{align}
    \text{FL}_{\text{{\textsc BLT}}} &= \text{Transf. FL}(h_{\mathcal{G}}, l_{\mathcal{G}}, m=n_{ctx}/n_p, V=0) / n_p\\
   &\quad +\text{Transf. FL}(h_{\mathcal{E}}, l_{\mathcal{E}}, m=w_{\mathcal{E}}, V=0) \\
   &\quad + \text{Transf. FL}(h_{\mathcal{D}}, l_{\mathcal{D}}, m=w_{\mathcal{D}}, V=256) \\ 
    &\quad + \text{Cross Attn. FL}(h_{\mathcal{E}}, l_{\mathcal{E}}, m=n_p, r=n_p/k) \cdot k/n_p \\ 
   &\quad + \text{Cross Attn. FL}(h_{\mathcal{D}}, l_{\mathcal{D}}, m=k, r=k/n_p) 
\end{align}

Here, $n_{ctx}$ denotes the length of the input sequence in bytes, $n_p$ specifies the patch size, and $r$ represents the proportion of queries relative to key/values. The variable $k$ indicates the ratio between the patch dimension and the byte dimension, effectively quantifying the number of individual model partitions that are combined to generate the overall model representation ($k=2$ is illustrated in \autoref{fig:crossattn}). The parameter $V$ is the vocabulary size utilized for output projection, which appears exclusively within the local decoder. The applicable attention context length, $m$, depends on whether a particular module operates over the byte sequence or the patch sequence. For each model component, we adjust the associated attention FLOPs to account for these differences. Detailed mathematical expressions for computing FLOPs in the Transformer and for Cross-Attention are outlined in Appendix~\ref{section:flops-appendix}.

\subsection{Bits-Per-Byte Estimation} Perplexity only makes sense in the context of a fixed tokenizer as it is a measure of the uncertainty for each token. When comparing byte and token-level models, following previous work~\citep{xue2022byt5, yu2023megabyte, wang2024mambabyte}, we instead report Bits-Per-Byte (BPB), a tokenizer independent version of perplexity. Specifically:

\begin{align}
    \text{BPB}(x)&= \frac{\mathcal{L}_{CE}(\pmb{x})}{\ln(2)\cdot n_{\text{bytes}}}
\end{align}

In this formulation, the aggregate cross-entropy loss representing the model’s uncertainty regarding the data $\pmb{x}$ is scaled by the total byte count of $\pmb{x}$ and a normalization factor.

\subsection{Transformer Architecture Hyperparameters} For each transformer block within {\textsc BLT}, encompassing both the local and global components, our approach closely mirrors the Llama~3~\citep{dubey2024llama} architecture. Specifically, feed-forward layers employ the SwiGLU activation function~\citep{swiglu}, while rotary positional embeddings (RoPE)~\citep{RoPe} with $\theta=500000$~\citep{xiong2024effective} are incorporated solely within the self-attention mechanisms. For normalization across layers, we utilize RMSNorm~\citep{RMSNorm}. All self-attention layers that apply fixed-standard attention patterns—such as \textit{block causal} or \textit{fixed-window block causal}—are implemented using Flash attention~\citep{dao2022flashattention}, where the window size for fixed-width masks is set to 512. For cross-attention modules, which require patch-dependent and thus dynamic masks, we adopt Flex Attention\footnote{\url{https://pytorch.org/blog/flexattention}} to facilitate fused operations and achieve substantial training speedups.

\subsection{{\textsc BLT}-Specific Hyperparameters} 

To assess the performance of {\textsc BLT} models, our experimental setup is two-fold: we examine both scaling behaviors and downstream task performance. Accordingly, we analyze models spanning a range of sizes—400M, 1B, 2B, 4B, and 8B parameters. Details of model architectures are documented in Appendix~Table~\ref{tab:arch}. 

For the first cross-attention layer within the local encoder, query vectors are initialized via max-pooling. Consistently across all {\textsc BLT} models, we employ $500,000$ hashes generated from a single hash function, setting n-gram sizes from 3 up to 8. All models are trained with a learning rate set at $4e-4$. This learning rate is selected based on hyperparameter sweeps over the interval $1e-3$ to $1e-4$ at the 400M and 1B scales, with results indicating that an identical rate is optimal for both token and {\textsc BLT} models. For scaling experiments with Llama-2 datasets, we adopt batch sizes as recommended by~\cite{dubey2024llama} or use the corresponding batch size in terms of bytes. Model optimization utilizes the AdamW optimizer~\citep{adamw} with hyperparameters $\beta_1 = 0.9$, $\beta_2 = 0.95$, and $\epsilon = 10^{-8}$. A linear learning rate warm-up is performed over 2000 steps, followed by cosine decay down to zero; weight decay is set at 0.1, and we apply global gradient clipping capped at 1.0.

\section{Scaling Trends}
\label{section:scaling}

We offer a comprehensive overview of the scaling behavior observed in byte-level models, providing insights that can guide the continued scaling of {\textsc BLT} models. Our scaling analysis addresses several shortcomings found in earlier investigations of byte-level modeling by: (a) contrasting scaling patterns under the regime of compute-optimal training, (b) training equivalent 8B parameter models with substantial datasets (reaching up to 1T tokens or 4T bytes) and assessing their effectiveness on downstream benchmarks, and (c) analyzing scaling performance under scenarios with controlled inference cost. In a subsequent section, we will explore the particular strengths that arise from modeling byte sequences directly.

\subsection{Parameter Matched Compute Optimal Scaling Trends}
\label{section:parameter-matched-scaling}
We leverage the Llama 2 dataset to train multiple \textit{compute-optimal} bpe and {\textsc BLT} models at four different scales, ranging from 1B to 8B parameters. Subsequently, we visualize the relationship between the number of training flops and the language modeling accuracy using a representative fraction of the training data mixture. The bpe models utilize an optimal balance of parameters and data, as specified in the Llama~3~\citep{dubey2024llama} recommendations. This \textit{compute-optimal} framework is theoretically expected to yield maximal training performance for the given compute resources~\citep{hoffmann2022training}, thus serving as a strong comparison point. Alongside each bpe model, an equivalent {\textsc BLT} model is trained under identical data conditions, with the Latent Transformer closely matching both the parameter count and architectural design of its bpe Transformer counterpart.

As shown in~\autoref{fig:scaling-compute-opt} (right), {\textsc BLT} models consistently match or surpass the performance of their bpe-based counterparts, and this pattern persists with increasing model scale and computation (flops). According to our knowledge, {\textsc BLT} is the first Transformer operating at the byte level to exhibit scaling behaviors comparable to BPE-based models in compute-efficient settings. This supports our hypothesis that the ideal proportion of parameters to training compute found in bpe-based architectures also holds for {\textsc BLT}, or is at least closely aligned.

Both enhancements in architecture and the introduction of dynamic patching are essential for aligning with bpe scaling trajectories. As shown in ~\autoref{fig:scaling-compute-opt} (left), we benchmark space-patching-based approaches against Llama 3. To estimate SpaceByte \citep{slagle2024spacebyte}, we use a {\textsc BLT} model with space-patching, omitting n-gram embeddings and cross-attention mechanisms. While SpaceByte demonstrates gains relative to Megabyte, its performance still lags significantly behind Llama 3. Meanwhile, \autoref{fig:scaling-compute-opt} (right) highlights the collective gains achieved via both refined architectures and dynamic patching techniques. At scale, {\textsc BLT} architectures reach comparable performance to leading tokenizer-based models like Llama 3.

Additionally, our findings indicate that, among models utilizing tokenizers, those employing the Llama-3 tokenizer achieve superior results compared to counterparts trained with the Llama-2 tokenizer, even when both are provided with identical training datasets.

In conclusion, the {\textsc BLT} architecture demonstrates scaling patterns that are intermediate between those of Llama 2 and Llama 3 when tested with substantially larger patch sizes. While the average token lengths for the bpe tokenizers in Llama 2 and 3 are 3.7 and 4.4 bytes, respectively, {\textsc BLT} attains analogous scaling behavior using mean patch sizes of 6 or even 8 bytes. Because inference flops decrease as the average patch size increases, selecting a patch size of 8 bytes results in nearly 50\% reduction in inference computation. Furthermore, models configured with larger patch sizes tend to yield superior performance as both model and data sizes scale up. Although {\textsc BLT} with an 8-byte patch size initially underperforms relative to the bpe variant of Llama 2 at the 1B scale, it surpasses bpe performance at the 7B parameter scale. These outcomes indicate that larger patch sizes could be advantageous for even greater model sizes, implying that expanding the patch size further might be possible as models and training computation budgets increase.

\subsection{Beyond Compute Optimal Task Evaluations}
\label{section:evals}
To further investigate scaling characteristics, we train an 8B {\textsc BLT} model at a parameter-to-token ratio that exceeds the compute-optimal limit, utilizing the {\textsc BLT}-1T dataset, which offers greater scale and improved data quality. The model's capabilities are then assessed across a range of widely adopted classification and generation benchmarks. For this purpose, we focus on benchmarks that test common sense reasoning, general knowledge, and code synthesis:
\paragraph{Classification tasks} comprise ARC-Easy (0-shot)~\citep{Eval_ARC}, Arc-Challenge (0-shot)~\citep{Eval_ARC}, HellaSwag (0-shot)~\citep{Eval_hellaswag}, PIQA (0-shot)~\citep{Eval_piqa}, and MMLU (5-shot)~\citep{hendrycks2020measuring}. We utilize a prompt-based likelihood estimation, scoring candidate responses by their predicted probabilities, and present the mean accuracy obtained.
\paragraph{Coding related generation tasks:} For code generation, we evaluate pass@1 performance on MBPP (3-shot)~\citep{Eval_MBPP} and HumanEval (0-shot)~\citep{Eval_HumanEval}, measuring the model's proficiency at generating correct Python code solutions.

In \autoref{tab:evals}, we present a comparative analysis of three models trained on the {\textsc BLT}-1T dataset: a bpe Llama 3 tokenizer-based model\footnote{The Llama 3 tokenizer, featuring a 128k vocabulary, was selected due to its superior performance over the 32k vocabulary in Llama 2.}, along with two versions of the {\textsc BLT} model. The two variants employ different patching strategies: {\textsc BLT}-Space leverages a space-patching approach, while {\textsc BLT}-Entropy uses a patching strategy guided by entropy. Both variants of the entropy-based model incorporate an approximately monotonicity-constrained method and restart the model context at each newline, as detailed in~\autoref{section:entropy-context}. All three systems are allocated the same flop budget for training. Notably, for {\textsc BLT}-Entropy, we decrease the entropy threshold during inference from 0.6 to 0.1, observing that this adjustment yields better task outcomes at the expense of requiring more inference steps.

The {\textsc BLT}-Entropy model achieves superior performance compared to the Llama 3 model on 4 of the 7 tasks, despite being trained on an equivalent byte count. This enhancement is probably attributable to (1) more efficient utilization of training computation through dynamic patching, and (2) the explicit modeling of byte-level representations rather than token-based ones.

Conversely, {\textsc BLT}-Space lags behind the Llama 3 tokenizer across nearly all tasks except one, though it does obtain substantial decreases in inference flops due to its notably larger mean patch size of 6 bytes. By contrast, the bpe and entropy-patching based approaches maintain a comparable mean patch size of around 4.5 bytes on the mixed training dataset. Under equivalent training resource constraints, the model employing the larger patch size is able to process 30\% more data than its counterparts, which could potentially drive {\textsc BLT} farther from the point of compute-optimality.

\subsection{Patches Scale Better Than Tokens} {\textsc BLT} models allow for the joint scaling of both model parameters and patch dimensions, all while preserving a fixed compute cost for training and inference and utilizing the same quantity of training samples. The flexibility to expand patch sizes without bound is distinctive to patch-based architectures; this capability enables them to circumvent the computational compromises inherent in token-based models with static vocabularies, as elaborated in Section~\ref{section:bpe-patch}. Utilizing longer patch sizes reduces computational requirements, permitting this saved compute to be redirected toward enlarging the global latent transformer, since it is invoked less frequently.

We perform a fixed inference scaling investigation to evaluate whether increasing model size, combined with fewer processing steps on larger patches, yields better performance compared to smaller models executing more steps. Using models of 400 million and 3.6 billion parameters initialized with the Llama 2 tokenizer, we identify flop-matched counterparts employing the Llama 3 tokenizer, as well as {\textsc BLT}-Entropy variants with mean patch sizes of 6 and 8 bytes within the training data mixture (refer to~\autoref{tab:fixed-inf-params} for specific model configurations). For the models utilizing patch size 8, we employ three encoder layers instead of one. Each configuration is trained under multiple total training flop constraints.

As depicted in \autoref{fig:fixed_inference_scaling}, {\textsc BLT} models demonstrate superior scaling behavior compared to architectures based on tokenization across both classes of inference FLOPs. Although BPE models tend to outperform at lower training compute levels, their advantages dissipate beyond the compute-optimal region, where {\textsc BLT} models begin to dominate. In practical applications, it can be advantageous to allocate greater resources during the one-off pretraining phase to secure enhanced performance under a fixed inference compute constraint. A notable instance of this strategy is represented by the 8B model category, such as Llama 3.1, which has been exposed to training datasets two orders of magnitude larger than the compute-optimal threshold typically recommended for models of its scale.

The point at which {\textsc BLT} surpasses token-based models has moved marginally closer to the compute-optimal threshold when transitioning to models with greater flop capacity—shifting from approximately 3x to 2.5x the compute-optimal expenditure. Likewise, the model with a larger patch size of 8 exhibits a more pronounced scaling behavior in the larger flop class, overtaking the other models at an earlier stage. As outlined in Section~\ref{section:parameter-matched-scaling}, scaling up patch sizes tends to yield performance approaching that of BPE models as overall model size increases. This trend can be partially explained by the reduced proportion of total flops devoted to the byte-level Encoder and Decoder modules, whose scaling is outpaced by that of the Latent Transformer. For example, expanding the total parameter count by 20 times—from 400M to 8B—results in just about a twofold increase in {\textsc BLT}'s parameters associated with its local model components. This detail is critical, given that enlarging patch sizes only affects the flop count of the patch Latent Transformer, leaving the byte-level modules unaffected. This explains why the ratio of {\textsc BLT}-Entropy with patch size 8 to the Llama 2 model size slightly increased from 1.6x to 1.7x when advancing to a larger scale model.

To conclude, our investigation into patch-length scaling reveals that the {\textsc BLT} patch-based framework benefits from joint enlargement of both patch and model sizes, leading to enhanced scaling behaviors. These positive effects appear to hold, or even strengthen, as model size continues to grow.

\section{Byte Modeling Improves Robustness} Additionally, we assess the resilience of {\textsc BLT} relative to models reliant on token-based representations without inherent byte-level access, and introduce a method for adapting pretrained token models to operate at the byte level.
\label{sec:robustness}

\subsection{Character-Level Tasks}

Initial interest in training byte-level models stemmed from their capacity to maintain stability in the presence of byte-level noise within inputs, as well as their inherent sensitivity to the internal structure of tokens—a challenge for traditional models reliant on tokenization. To systematically assess these characteristics, we conduct supplementary tests using benchmarks designed to evaluate resilience against input corruption and recognition of fine-grained elements, such as individual characters in English and multiple languages, along with digits and phonemes. The findings from these evaluations are reported in Table~\ref{tab:char_tasks}.

\paragraph{Noisy Data} To assess the robustness of {\textsc BLT} relative to tokenizer-based approaches, we generate corrupted versions of the benchmark classification datasets discussed in Section~\ref{section:evals}. Five character-level noise techniques are used to systematically alter the text: (a) \textit{AntSpeak}: The entire input is rewritten as uppercase letters, separated by spaces. (b) \textit{Drop}: Randomly deletes 10\% of the text’s characters. (c) \textit{RandomCase}: Randomly assigns 50\% of characters to uppercase and the remaining 50\% to lowercase. (d) \textit{Repeat}: Selects 20\% of characters and replicates them up to four additional times. (e) \textit{UpperCase}: Converts every character in the input to uppercase. For each evaluation, every noise method is independently applied to the prompt, the completion, or both, and performance metrics are averaged across these variations. Table~\ref{tab:char_tasks} summarizes the outcomes for HellaSwag~\citep{Eval_hellaswag} under noisy conditions, showing that {\textsc BLT} consistently exhibits greater resilience than tokenizer-based baselines—averaging an 8-point gain over models trained on identical corpora—and even surpasses the Llama 3.1 model, which was trained using substantially more data.

\paragraph{Phonology - Grapheme-to-Phoneme (G2P)} We evaluate the proficiency of {\textsc BLT} in converting sequences of graphemes (alphabetic symbols denoting a word) into their corresponding phonemic transcriptions (representations of pronunciation). The outcomes for the G2P task, conducted within a 5-shot framework and utilizing Phonology Bench~\citep{suvarna2024phonologybench}, are displayed in Table \ref{tab:char_tasks}. These results demonstrate that {\textsc BLT} achieves superior performance compared to the Llama 3 1T model employing a tokenizer-based approach.

\paragraph{CUTE}
To measure character-level comprehension, we test {\textsc BLT} using the CUTE benchmark~\citep{edman2024cute}, which contains a variety of tasks organized into three main types: compositional understanding, orthographic similarity, and sequence manipulation. The CUTE tasks present notable difficulties for typical tokenizer-based models, as they tend to retain awareness of how tokens are spelled, yet often fall short when asked to exploit this for text manipulation. Results in Table~\ref{tab:char_tasks} reveal that {\textsc BLT}-Entropy surpasses both BPE Llama 3 baselines by over 25 points on the benchmark as a whole. Notably, our model excels on the character manipulation subcategory, obtaining 99.9\% accuracy for each spelling-oriented task. These considerable gains, despite {\textsc BLT} being trained with only one-sixteenth the data volume of Llama 3.1, suggest that acquiring robust character-level representations is challenging for BPE models. Figure \ref{fig:cute_examples} provides several examples where the Llama 3 tokenizer model falters while our {\textsc BLT} model is successful. Among the tasks, BPE outperforms only on word deletion and word insertion, perhaps reflecting the relative difficulty for byte-level models to perform word-scale manipulations; however, the difference in performance on these is not substantial, and it is arguably easier to assemble words from characters than to decompose words within BPE frameworks. Our evaluation methodology maintains consistency across all tasks and employs the original Huggingface prompts. It is possible that further refinements in prompt design could enhance BPE model outcomes.

\paragraph{Low Resource Machine Translation} We assess the performance of {\textsc BLT} on translation tasks involving twenty-one lower-resource languages representing six major language families, each with diverse scripts, using the FLORES-101 benchmark~\citep{goyal2022flores}. SentencePiece BLEU scores, detailed in Table~\ref{tab:flores}, are used for evaluation.
The findings indicate that {\textsc BLT} achieves a 2-point overall BLEU increase when translating into English and is 0.5 points better when translating from English compared to a counterpart trained with the Llama 3 tokenizer. For widely spoken language pairs, {\textsc BLT} matches or slightly surpasses Llama 3's translation quality. Notably, in many lower-resource language family pairs, {\textsc BLT} demonstrates clear superiority over Llama 3, highlighting the advantages of byte-level modeling in effectively generalizing to underrepresented, long-tail byte sequences.

\subsection{Training {\textsc BLT} from Llama 3}
\label{sec:distillation}
We investigate a strategy in which {\textsc BLT} models benefit from established pre-trained tokenizer-based models to promote more rapid and robust convergence during training. This approach initializes the global transformer parameters in {\textsc BLT} using those from a pre-trained Llama 3.1 model. Following this initialization, we fine-tune the global transformer parameters with a learning rate set to one-tenth of that applied to the local encoder and local decoder components, matching the optimal training duration for Llama 3. Our results, detailed in Table~\ref{tab:distill}, allow direct comparison with a baseline {\textsc BLT}. The findings clearly show that {\textsc BLT} derived from Llama 3.1 substantially outperforms both the original Llama 3 and baseline {\textsc BLT} models when provided with an equal computational budget in terms of flops. Notably, when evaluated against the {\textsc BLT}-Entropy variant (refer to Table~\ref{tab:evals}), which underwent training on a much larger corpus (1T tokens), the {\textsc BLT} variant initialized from Llama 3.1 still attains better results on the MMLU benchmark. This indicates the proposed workflow offers a promising method to notably diminish the training flops while maintaining or improving performance.

This approach can be interpreted as adapting tokenizer-based models to operate without tokenizers, thereby transforming a pre-trained LLaMA 3.1 model into a {\textsc BLT} model. For a thorough evaluation, we report the results of the original LLaMA 3.1 trained on 15T tokens in Table \ref{tab:distill} and assess its performance relative to the {\textsc BLT} variant derived from LLaMA 3. Although our model shows only a modest decrease in MMLU and HumanEval scores, it undergoes a more pronounced performance reduction on other benchmarks. These findings highlight the need for continued research to maximize the benefits of pre-trained models, including refining the composition of data mixtures and tuning relevant hyperparameters.

\section{Ablations and Discussion}
\label{section:ablations}
Here, we present a series of ablation studies that support the architectural decisions underlying {\textsc BLT}. Additionally, we analyze the patching approach and the selection of hyper-parameters employed for the {\textsc BLT} model with 8B parameters, which was trained on the {\textsc BLT}-1T dataset.

\paragraph{Entropy Model Hyper-parameters} To investigate how changes in entropy model scale and context window size affect performance scaling, we conduct experiments training byte-level entropy transformer architectures ranging from 1 million to 100 million parameters, each with context windows spanning 64 to 512 tokens. We illustrate the relationship between bits-per-byte (bpb) and training FLOPs by plotting scaling law curves, utilizing results from our $400m$ and $1b$ {\textsc BLT} models that have been trained on the Llama-2 dataset (see \autoref{fig:entropy_ablation}). Our analysis reveals that increasing both the size of the entropy model and its context window improves scaling performance; however, gains diminish when model size exceeds 50 million parameters.

\paragraph{Types of Patching}

We perform an ablation study on the four patching strategies detailed in Section~\ref{section:patching}: (1) Strided Patching utilizing strides of 4 and 6, (2) Whitespace-Based Patching, (3) Patching according to Byte Pair Encoding Tokenizer derived from the Llama 3 tokenizer, and (4) Entropy-Based Patching, which leverages a compact byte-level LLM.

Although dynamic patching decreases the actual sequence length, we ensure the context length remains comparable across all patching strategies by regulating sequence length. Each model is exposed to an equivalent number of bytes per sequence on average during both training and inference, thus avoiding any biases from differing context sizes. The outcomes of these ablation studies are summarized in Figure~\ref{fig:scaling-compute-opt}. Notably, all non-static patching approaches demonstrate superior performance relative to static patching, and space patching yields results that are nearly indistinguishable from those obtained with dynamic entropy based patching.

In \autoref{tab:evals_ablation}, we report benchmark results for {\textsc BLT} models, offering a comparison between approaches that utilize tokenizers, space patching, and entropy-based patching, all trained on the Llama 2 dataset for a tuned number of steps~\citep{dubey2024llama}. While space patching constitutes a more straightforward method—eschewing the computational overhead of computing entropies during training—we observe that the advantages yielded by entropy-based patching in scaling analyses~(Section~\ref{section:scaling}) consistently persist and extend to performance on downstream benchmarks.\footnote{Results from space patching correspond to earlier experiments that did not use cross-attention; nevertheless, comparable patterns are present in settings with cross-attention.}

\paragraph{Cross-Attention} 
In~\autoref{tab:crossattn_ablations_1b}, we investigate the effects of introducing cross-attention at different locations within both the encoder and decoder components of {\textsc BLT} through an ablation study. For the encoder’s cross-attention mechanism, we experiment with three initialization strategies for the queries: (1) using a single shared learned embedding across all global states, (2) applying a hash-based embedding to represent the bytes within each patch, and (3) computing a pooled summary from the encoder’s hidden states corresponding to the patch bytes at a particular encoder layer.

Our results indicate that integrating cross-attention within the \textit{decoder} yields the highest effectiveness. Incorporating cross-attention into the encoder leads to minor gains, primarily when queries are initialized via pooling. Moreover, the advantages of cross-attention are more pronounced on the Common-Crawl dataset and become especially significant as patch sizes increase.

\paragraph{n-gram Hash Embeddings} 

We conduct an ablation study with n-gram hash embedding vocabulary sizes of 0, 100K, 200K, and 400K, with the corresponding results shown in Table~\ref{tab:ngram_ablations_1b}. Our analysis reveals that incorporating hash embeddings leads to improvements across all tested domains, with the most pronounced impact observed on Wikipedia and Github, where the improvement reaches 0.04 bpb compared to just 0.01 bpb after 15k training steps at the 8B parameter scale. When reducing the number of hashes from 500K to 300K at the 8B scale, the resultant change in performance is minimal—around 0.001 bpb—after 15k training steps. These findings suggest that hash embeddings are essential for {\textsc BLT} models to achieve comparable results to those of tokenizer-based approaches; nevertheless, performance gains begin to plateau beyond 300K hashes. Furthermore, our observations suggest that the benefits of hash embeddings largely complement those obtained from cross-attention, as each contributes improvements in distinct datasets.

\paragraph{Local Model Hyperparameters} Table \ref{tab:local_layers} presents an ablation study exploring different layer configurations for both the local encoder and decoder. Notably, when using hash n-gram embeddings, {\textsc BLT} achieves strong results with a highly minimal encoder, consisting of only a single layer, while relying on a comparatively more complex, deeper decoder.

\section{Related Work}
\label{section:related_work}

\textbf{Character-Level RNNs:} The task of Character Language Modeling has remained central in neural model research since its inception~\citep{sutskever2011generating,mikolov2012subword,graves2013generating}, primarily because these models can naturally handle previously unseen words, bypassing traditional back-off strategies. \cite{kim2016character} introduced a system employing convolutional and highway layers to process input at the character level, with their outputs serving as input to LSTM-based RNNs. This approach achieved parity with leading RNN-based language models of that era for English, and even surpassed them in morphologically complex languages—a key motivation for adopting character-level language modeling. In a related effort, \cite{kenter2018byte} demonstrated that byte-level LSTM models outstripped the performance of word-level alternatives for machine comprehension tasks on morphologically intricate languages such as Turkish and Russian. Similarly, \cite{zhang2015character} deployed character-based convolutional neural networks for classification, observing that these outperformed their word-based counterparts on particular benchmarks. Additionally, \citet{chung2019hierarchical} proposed hierarchical LSTM architectures equipped with boundary detectors at multiple levels to identify implicit hierarchical structures within text, which further enhanced the efficacy of character-level language models. Moreover, ByteNet~\citep{kalchbrenner2016neural} replaced attention mechanisms with convolutional layers operating on characters to tackle machine translation.

\textbf{Character-Level Transformers:} The introduction of attention-based transformer architectures~\citep{vaswani2017attention} and the use of subword tokenization techniques~\citep{sennrich-etal-2016-neural} have notably enhanced the effectiveness of neural approaches in tasks such as language modeling and across a variety of benchmarks. Nonetheless, utilizing word or sub-word representations inherently sets a preference for the model's abstraction level. In an effort to merge the advancements of transformer methodologies with earlier encouraging findings in character-based language modeling, \citet{al2019character} implemented exceptionally deep transformers. Supported by auxiliary loss functions, they trained transformer models that surpassed the performance of prior LSTM-based character language models. Still, a considerable performance discrepancy persisted in favor of word-level LLMs. Additionally, GPT-2~\citep{radfordgpt2} documented that, for large-scale datasets such as the 1 Billion Word Benchmark, language models trained at the byte level failed to match the outcomes achieved by their word-level counterparts.

Although \cite{Choe2019BridgingTG} showed that transformer-based LLMs operating at the byte level can surpass subword-level LLMs when controlling for parameter count, such models demand significantly greater computational resources and longer training times. Likewise, \cite{el2020characterbert} introduced CharFormer, a BERT variant that constructs word representations by applying convolutional layers to character embeddings, reporting improved outcomes particularly in the medical domain—but at the cost of substantially higher compute usage. In the same vein, \cite{clark2022canine} proposed CANINE, an encoder-only model with 150M parameters designed to process input directly at the character level. At its core, CANINE leverages a deep transformer stack reminiscent of our global model, complemented by a local transformer and strided convolution operations to reduce the character sequence, outperforming mBERT (a token-level encoder-only model) on several downstream multilingual benchmarks. ByT5~\citep{xue2022byt5} examined byte-level encoder-decoder architectures without integrating any patching mechanisms; their model achieved increased robustness to noise and matched the performance of tokenizer-based models using only a quarter of the data. However, because these models omitted patching, they required computing attention across the full byte sequences, resulting in substantial computational burden. Using bytes rather than subword units leads to considerably longer input sequences, posing challenges for scaling byte-level architectures efficiently. More recently, \cite{wang2024mambabyte} explored byte-level modeling via the Mamba Architecture~\citep{gu2023mamba}, which allows for fixed-size memory states over extensive contexts. Without any patching, their byte-level Mamba model demonstrated improved efficiency, outperforming similarly-sized byte-level transformers (at 350M parameters) with respect to bits-per-byte across several datasets in a compute-matched (FLOP-controlled) comparison.

\textbf{Patching-based approaches:} Leveraging patching methods can significantly reduce the excessive number of floating point operations typically required by byte-level LLMs, potentially maintaining model efficacy. Several studies have reported promising results on smaller model sizes and limited training datasets. For example, \cite{nawrot-etal-2022-hierarchical} examine static patching techniques for downsampling and upsampling, introducing the hourglass transformer, which surpasses other byte-level models at the 150M parameter level. Building on this, \cite{nawrot-etal-2023-efficient} enhance performance through dynamic patching strategies. These include a learned boundary-predictor trained end-to-end, a boundary-predictor guided by selected tokenizers, and an entropy-driven patching framework resembling {\textsc BLT}. They demonstrate that such models outperform standard transformers of their period on language modeling tasks using roughly 40M parameters and 400M training tokens. Additionally, \cite{lester2024training} explore training on sequences that have been compressed using arithmetic coding, achieving compression rates superior to those obtained by BPE. Their use of equal-info windows enables surpassing byte-level benchmarks in language modeling tasks, although subword-level baselines still show superior performance.

Our study is heavily informed by MegaByte~\citep{yu2023megabyte}, a decoder-only, causal large language model that leverages a fixed, static patching approach combined with representation concatenation to segment bytes into patches. On the decoding side, MegaByte employs a local model to reconstruct the original byte sequences from these patches. The authors of MegaByte demonstrate that their model achieves performance on par with models using tokenizers when evaluated at the 1B parameter scale on data comprising approximately 400B bytes. In our experiments, we provide ablation studies of MegaByte and consistently observe that its static patching strategy falls short of matching the performance of compute-optimal, state-of-the-art tokenizer-based models within a controlled FLOP regime; we further illustrate how {\textsc BLT} narrows this performance gap. Similarly, \cite{slagle2024spacebyte} report this performance disparity for MegaByte and propose enhancements by extending static patching to operate on whitespaces and similar byte classes, as well as incorporating a local encoder module. Their modifications lead to performance gains relative to tokenizer-based transformers in certain compute-constrained scenarios, notably in domains such as Github and arXiv at the 1B parameter scale. We replicate and report our own findings with this variant, concluding that even with such modifications, further refinements to the underlying model architecture are essential for byte-level models to truly scale and to approach or exceed the performance levels of leading token-based systems like Llama 3.

\section{Limitations and Future Work}

For the design of our models, we adhere to the number of training steps found to be optimal for Llama~3~\citep{dubey2024llama}. Nevertheless, it is important to note that the original scaling laws were established with BPE-level transformers in mind and may not yield the ideal ratio of data to model parameters when applied to {\textsc BLT}. Deriving scaling laws specifically tailored for {\textsc BLT}, which may further enhance scalability, is a direction we reserve for future research. Furthermore, much of our experimental work has been limited to models with up to 1B parameters; it remains plausible that, as we extend to models with 8B parameters or more, the preferred architectural configurations may evolve, potentially resulting in superior outcomes for larger-scale models.

Most current transformer frameworks and codebases are optimized for maximal efficiency when used with tokenizer-based transformer models. Although our study includes theoretically flop-matched comparisons and leverages advanced techniques like FlexAttention to efficiently process layers that differ from the standard transformer design, our current implementations might still lag behind tokenizer-based architectures regarding wall-clock runtime and could be improved with additional optimization efforts.

Although {\textsc BLT} employs an entropy model for patching that is trained separately, exploring the possibility of jointly learning the patching model in an end-to-end manner represents a promising avenue for further investigation. Section \ref{sec:distillation} provides preliminary experiments that suggest successful adaptation of tokenizer-based models, such as Llama 3—which have been trained on over 10T tokens—into byte-level models. This is achieved by initializing the global transformer with pretrained weights and subsequently freezing them. Advancing this line of research could yield techniques that preserve the advantages of byte-level representations while also potentially surpassing the performance of the original tokenizer-based models, without necessitating complete retraining.

\section{Conclusion}
\label{section:conclusion}

In this work, we introduce the Byte Latent Transformer (\textbf{BLT}), an innovative architecture that departs from the standard reliance on predefined token vocabularies in large language models. {\textsc BLT} employs a novel, learnable mechanism that adaptively assembles bytes into variable-length patches, thereby tailoring computational effort in accordance with the intricacy of the input data. This approach translates to notable gains in both computational efficiency and model robustness. Through comprehensive scaling experiments, we show that {\textsc BLT} achieves performance on par with established token-based systems such as Llama 3 across scales of up to 8B parameters and 4T processed bytes, while also enabling up to 50\% savings in inference flops for only slight sacrifices in standard evaluation metrics. In addition, {\textsc BLT} facilitates simultaneous scaling of both model size and patch granularity without breaching a set inference compute threshold—an advantage that is particularly relevant for typical compute-limited scenarios. By operating directly at the byte level, {\textsc BLT} enhances model resilience to atypical and noisy data, offering stronger handling of the data distribution’s long-tail and greater insight into subword-level regularities. Taken together, these advances position {\textsc BLT} as a compelling substitute for conventional tokenization methodologies, combining scalability, adaptability, and robustness within a unified language modeling framework.

\end{document}